% ============================================================
\section{Conclusions}
\label{sec:conclusions}
% ============================================================

We have presented the \emph{CRNS Site Characterization Pipeline},
a fully integrated Python tool for automating the pre-installation
assessment of Cosmic-Ray Neutron Sensing stations from global open
datasets.
The pipeline addresses a recognised gap in the CRNS ecosystem:
while operational correction routines are available (e.g.\
\texttt{crnpy}), no single tool previously computed all
\emph{site-specific} correction factors and supplementary
characterisation quantities before sensor deployment.

\paragraph{Core physics outputs.}
The three multiplicative correction factors — $\kT$ (topographic,
3-D ray-casting), $\kM$ (muon field-of-view, analytical integral),
and $\kL$ ($W(r)$-weighted WorldCover + OSM) — are combined into
$\kTot = \kT\kM\kL$.
Together with the lattice water $\lw$ derived from SoilGrids clay
fractions, they determine the reference count $N_0$ through the
Desilets (2010) formula, which preserves full analytical
invertibility and remains the primary calibration reference.

\paragraph{Supplementary characterisation.}
The seasonal climatology of atmospheric water vapour ($f_{WV}$),
snow-water equivalent (SWE), and above-ground biomass hydrogen
(AGBH) quantifies the operational corrections that will be required
throughout the measurement period.
The Macrostrat geological characterisation flags unusual lithologies
that may bias the clay-based $\lw$ estimate.
The Topographic Wetness Index provides a spatially distributed map
of expected moisture variability within the footprint, relevant to
calibration planning.

\paragraph{Sampling design.}
The $W(r)$-optimised 26-point sampling plan, and the accompanying
$\rzz(\thetav)$ dynamic-range plot, translate the footprint physics
directly into field-campaign instructions.
The plan minimises the variance of the footprint-averaged soil
moisture estimate for a fixed number of samples.

\paragraph{Uncertainty.}
A combined uncertainty budget for the Italian Alps example yields
$\sigma_{N_0}/N_0 \approx 5\%$, with the LULC classification
uncertainty ($\sim 4\%$) and the soil lattice water ($\sim 3\%$)
as the dominant terms.
This demonstrates that accurate site characterisation can
substantially reduce the systematic error in soil-moisture
retrievals, and identifies where additional field effort (in-situ
clay measurements, vegetation water content surveys) is most
rewarding.

\paragraph{Outlook.}
Future extensions could include:
\begin{itemize}[nosep]
  \item Integration of LiDAR DEMs where available, to resolve
        sub-30-m terrain features relevant for $\kT$.
  \item A sensitivity analysis module that propagates input
        uncertainties through the full correction chain and reports
        quantitative bounds on $N_0$.
  \item Time-varying $\kL$ driven by satellite-derived LAI and
        canopy water content, enabling quasi-operational correction.
  \item Support for portable CRNS rovers \citep{schron2017}, where
        the footprint geometry changes with vehicle speed and the
        sampling plan module must be adapted accordingly.
\end{itemize}

The pipeline source code is openly available at
\url{https://github.com/barbierimauro/site_characterization}
under an open-source licence.
All datasets used are publicly accessible without registration.
We invite contributions and reports of site-specific validation
datasets to expand the empirical basis for $f_H$ assignments and
$\lw$ estimates.
